\documentclass[../main.tex]{subfiles}
\begin{document}

\section{Introduction}
\label{sec:intro}

Graph neural networks (GNNs) have become the default architecture for social bot
detection. By treating accounts as nodes and interactions as edges, models such as
BotRGCN~\cite{feng2021botrgcn} and its successors exploit the observation that
automated accounts are easier to identify collectively than individually. The
premise underlying this entire family of methods is \emph{homophily}: the
assumption that an account's neighbours are informative about its label, so that
averaging over a neighbourhood sharpens rather than blurs the signal.

Modern bots attack exactly this premise. Rather than forming the dense, mutually
connected ``bot-nets'' of earlier generations, contemporary accounts engage in
\emph{relation camouflage}~\cite{dou2020caregnn,cresci2020decade}: they
deliberately follow and interact with large numbers of genuine users, so that a
bot's neighbourhood looks statistically like a human's. Under standard message
passing, the bot's anomalous features are averaged against benign neighbours and
the discriminative signal is diluted---a failure mode that recent bot-detection
work reports consistently~\cite{he2025bothp,dang2026xhbot}.

The empirical response has been architectural. Detectors now prune suspicious
edges, split aggregation into homophilic and heterophilic channels, or disentangle
behavioural from positional signals; XHBot~\cite{dang2026xhbot} combines all three.
These designs demonstrably help, but the evidence for \emph{why} they help is
almost entirely ablation-driven. Three practical questions therefore remain
unanswered:

\begin{enumerate}
    \item \textbf{When is a graph-based detector worth deploying at all?} If
    camouflage is severe enough, averaging over a neighbourhood must eventually do
    more harm than good. Practitioners have no criterion for where that point lies.
    \item \textbf{What does edge pruning actually buy?} Pruning modules are
    justified by ablation deltas, not by a statement of the form ``this module
    extends the operating regime by \emph{this much}''.
    \item \textbf{Are multi-channel aggregators genuinely more expressive, or
    merely better tuned?} An architecture with more parameters that wins an
    ablation has not thereby been shown to be more capable.
\end{enumerate}

This paper answers these questions analytically. We adopt a contextual stochastic
block model (CSBM)~\cite{deshpande2018contextual} augmented with an explicit
adversarial parameter---the \emph{camouflage ratio} $\rho$, the fraction of a bot's
edges directed at humans---and study what one round of mean aggregation does to
class separability as $\rho$ grows. This places the work in the line of analysis
opened by Baranwal et al.~\cite{baranwal2021graph}, who showed that graph
convolution improves linear separability by a factor of roughly $1/\sqrt{D}$ for
neighbourhood size $D$ under a homophilic generative process.

\paragraph{What is already known.} We are explicit about priority from the outset.
For a \emph{balanced} symmetric CSBM, Ma et al.~\cite[Thm.~2]{ma2022homophily}
already established a closed-form threshold at which aggregation ceases to beat the
graph-free classifier, and our critical ratio reduces to their condition exactly when
the two classes are equal in size (Remark~\ref{rem:ma}); related crossover analyses
appear in~\cite{luan2023when,mao2023demystifying,fountoulakis2023graph}. The
existence of a crossover is therefore not our discovery, and we do not claim it.
Our contribution is the \emph{class-imbalanced adversarial} case, where the two
classes deflate by different amounts---which moves the decision boundary and makes
bot prevalence itself a robustness parameter---together with a closed form for how
far edge pruning shifts the threshold.

\paragraph{Contributions.}

\begin{itemize}
    \item \textbf{The edge-balance identity, and bot prevalence as a robustness
    parameter.} Camouflage edges have two ends, so a bot that points an edge at a
    human necessarily pollutes that human's neighbourhood too. We show this forces
    $\eta = c\rho$, where $c$ is the bot-to-human ratio, so that the total deflation
    is $1-(1+c)\rho$ (Proposition~\ref{prop:deflation}). The consequence is that
    bot prevalence enters the robustness threshold directly: as bots become more
    common, each one needs \emph{less} camouflage to break the detector. This is not
    a class-imbalance problem that resampling addresses.

    \item \textbf{A Bayes-error crossover valid under class imbalance.} When
    $c \neq 1$ the two classes deflate asymmetrically, the post-aggregation class
    midpoint moves, and a comparison against a fixed decision boundary is no longer
    valid---the objection Luan et al.~\cite{luan2023when} raise
    against~\cite{ma2022homophily}. We compare Bayes errors instead
    (Theorem~\ref{thm:error}), obtaining
    $\rho^\star = (1 - D^{-1/2})/(1+c)$ (Corollary~\ref{cor:critical}), which
    recovers~\cite{ma2022homophily} at $c=1$ and~\cite{baranwal2021graph} at
    $\rho=0$. Our error term also retains a mixture-variance contribution that a
    naive $\sigma^2/D$ calculation misses, and which shifts the threshold by $13\%$.

    \item \textbf{A closed form for what edge pruning buys.} We show the
    Dirichlet-energy score underlying SGTR separates camouflage from genuine edges
    by exactly $\|\Delta\|^2$ in expectation (Lemma~\ref{lem:energy}), and prove
    that soft pruning raises the tolerated camouflage to
    $\rho^\star_{\mathrm{SGTR}} = \rho^\star / (\rho^\star + (1-\bar e_c)(1-\rho^\star))$
    for a pruner of quality $\bar e_c$ (Theorem~\ref{thm:sgtr}). The underlying
    reweighting identity is known~\cite{weiss2026cost,fountoulakis2023graph}; what we
    add is its composition with the finite-$D$ threshold to bound the
    \emph{adversary's} budget, which requires a regime those analyses do not work in
    (Remark~\ref{rem:sgtr-priority}). Because $\bar e_c$ is measurable for a deployed
    pruner, this converts an ablation delta into an operating-range claim, and makes
    explicit that no imperfect pruner attains immunity.

    \item \textbf{An honest account of the disentanglement objective.} We show the
    contrastive prototype loss maximises a variational lower bound on the mutual
    information between behavioural code and label
    (Proposition~\ref{prop:nce}), and prove the orthogonality penalty controls
    code dependence \emph{only} under a joint-Gaussian assumption
    (Proposition~\ref{prop:mi}). We state plainly where this falls short of
    identifiability, in light of known impossibility results for unsupervised
    disentanglement~\cite{locatello2019challenging}.

    \item \textbf{A ceiling on what pruning can achieve.} Down-weighting edges
    non-uniformly also forfeits effective degree, so pruning trades bias against
    variance rather than buying the former for free
    (Proposition~\ref{prop:deff}). Accounting for this corrects a conclusion that
    the composition analysis alone suggests: even a \emph{perfect} pruner tolerates
    only $\rho < 1 - 1/D$, not $\rho \to 1$ (Corollary~\ref{cor:ceiling}). We then
    bound the achievable pruning quality from the exact law of the edge score,
    which removes the free parameter: the tolerance becomes a function of
    $(D, c, r, d)$ alone.

    \item \textbf{Depth, and an upper edge to the harmful region.} We give the
    exact mean--variance recursion across layers (Theorem~\ref{thm:depth}) and show
    the crossover is \emph{independent of depth}, while the gain or loss compounds
    exponentially in it. Because only $|\lambda|$ matters, the harmful region is a
    bounded \emph{window}: beyond it a bot's neighbourhood is strongly
    anti-correlated with its label and so informative again. The dangerous strategy
    is intermediate camouflage, not maximal camouflage.

    \item \textbf{Measurement on real graphs, and validation.} Every closed form is
    checked against Monte-Carlo simulation. Beyond that we \emph{measure} $\rho$ on
    two real labelled fraud graphs, finding $\hat\rho \approx 0.80$--$0.90$, just
    inside the predicted harmful window, and confirming the edge-balance identity
    to within $3\%$ where its degree assumption holds
    (Section~\ref{sec:measuring}). We are explicit that these are fraud rather than
    bot graphs, and that access to TwiBot was unavailable.
\end{itemize}

Our aim is not to propose another detector. It is to supply the missing analytical
account of a design pattern that the field has already converged on empirically,
and in doing so to replace a heuristic with a criterion that a practitioner can
check.

\end{document}
